\chapter{Metodologia}

Este capítulo descreve a metodologia empregada para o desenvolvimento e implementação do modelo preditivo de preços de ações, incluindo as etapas de coleta e pré-processamento dos dados, desenvolvimento e treinamento dos modelos, e a avaliação de seu desempenho. A metodologia adotada visa garantir a precisão e a confiabilidade das previsões, abordando cada fase do processo de forma estruturada e sistemática.

\section{Coleta de dados}

A fase inicial do processo envolve a coleta de dados, com o objetivo de obter informações detalhadas sobre o histórico de preços, volume de negociações, indicadores de análise técnica e outros fatores relevantes. Para isso, foram exploradas fontes como o site da B3, ligado à bolsa de valores de São Paulo (Bovespa), que oferece uma ampla gama de indicadores. Além disso, foi utilizada a API gratuita do site MarketStack, que disponibiliza informações sobre o valor de uma determinada ação em um intervalo de período de tempo que varia entre 10 minutos e cinco horas. Optei pelo intervalo de tempo de 15 minutos, pois, desta forma, a API fornece quatro \textit{candles} por hora, o que permite obter mais dados sobre o comportamento da ação ao longo do dia.

No âmbito deste trabalho, foi acessada a API para obter os valores da ação AAPL da Apple em intervalos de tempo de 30 minutos, no período compreendido entre as 00:00 horas do dia 19 de Julho de 2019 e as 19:45 horas do dia 14 de Junho de 2023. Esses dados serão fundamentais para análises posteriores e para a tomada de decisões no contexto do estudo, escolhi este intervalo de tempo pela quantidade de informações que este período oferece, gerando um arquivo csv com 35.126 \textit{candles}, optei por coletar os dados em 4 anos porque com isso, os modelos levariam em consideração um contexto mais recente do que em um intervalo de tempo maior. Gerei um arquivo csv chamado "AAPL\_data.csv" onde coloquei os resultados da consulta á API, esse arquivo contou com 35.126 registros de \textit{candles}, possuindo um tamanho de 1984 KB.

\section{Pré-processamento dos dados}

Após a coleta, o trabalho segue para a etapa do pré-processamento dos dados coletados, uma atividade crítica para garantir a qualidade dos dados. Nesta fase, são eliminados valores indesejados, como preços negativos ou nulos, visando assegurar a integridade e confiabilidade dos dados. Todos os valores resultantes são verificados para garantir conformidade com as normas do mercado financeiro.

Os indicadores de preços utilizados para cada intervalo de tempo, ou “candle”, incluem o valor da ação na abertura do intervalo (\textit{Open}), o maior valor atingido durante o intervalo (\textit{High}), o menor valor atingido durante o intervalo (\textit{Low}), e o valor ao fechamento do intervalo (\textit{Close}). Esses quatro valores, coletivamente conhecidos como OHLC (Open, High, Low, Close), são essenciais para a análise técnica e para a modelagem de séries temporais.

A API apresentou dados em um formato compatível com o era esperado para utilização no presente trabalho, com dados dos preços sendo do tipo ponto flutuante com 2 ou 3 casas decimais. Por conta disso, não houve necessidade de realizar tratamentos adicionais destes dados. Para este trabalho, o valor \textit{Open} foi escolhido como o principal parâmetro de treinamento da rede neural, devido ao seu menor índice de dados nulos comparado aos outros indicadores. Dados como \textit{Close} frequentemente apresentam valores nulos nas respostas da API, o que pode comprometer a qualidade das previsões. Portanto, a escolha do valor \textit{Open} visa garantir uma base de dados mais completa e confiável para o treinamento do modelo preditivo, cada valor \textit{Open} se refere ao \textit{candle} específico e não ao dia.

\section{Desenvolvimento}

Com os dados coletados e normalizados, inicia-se a etapa seguinte, que se concentra no desenvolvimento, treinamento e validação das redes neurais que modelam o problema. Para essa etapa, foi escolhida a plataforma \textit{Google Colab}, um produto gratuito da Google baseado no Projeto Jupyter. Essa escolha se justifica pelo poder computacional fornecido pela nuvem da Google, permitindo a execução eficiente de grandes conjuntos de dados. O Google Colab utiliza a linguagem Python e, para o desenvolvimento neste projeto, optou-se pelo framework Keras.

Durante o desenvolvimento, diversos modelos de redes neurais foram implementados. Cada modelo foi treinado e avaliado separadamente para identificar o que melhor se ajusta aos dados e objetivos do projeto. Resultados são reportados e os códigos que os geram são guardados para facilitar a reprodutibilidade e comparação entre os diversos modelos treinados. Este processo de comparação permite ajustar hiperparâmetros e otimizar a arquitetura dos modelos com base em métricas de desempenho específicas para cada tipo de treinamento, garantindo que o modelo final escolhido seja o mais adequado para a previsão de preços de ações.

Durante a fase de treinamento, o algoritmo é capacitado a aprender sobre os dados e seus padrões, enquanto o conjunto de validação é usado para comparar o desempenho do algoritmo com dados não presentes no conjunto de treinamento.

\section{Modelos}

Os modelos escolhidos para o trabalho foram o LSTM (Long Short-Term Memory), GRU (Gated Recurrent Unit) e TCN (Temporal Convolutional Network), todos reconhecidos por sua eficácia em problemas de séries temporais, onde dados são apresentados ao longo de um intervalo de tempo. Esses modelos são particularmente adequados para a previsão de valores de ações, que exibem comportamentos variados e complexos ao longo do tempo. A escolha desses modelos se baseia no referencial teórico, que demonstra que LSTM, GRU e TCN são eficazes em tarefas semelhantes, proporcionando resultados significativos na previsão de séries temporais e na captura de padrões temporais nos dados financeiros.

\section{Avaliação dos Modelos}

Os critérios para a avaliação do melhor modelo passam pelo comportamento de cada modelo durante a etapa de treinamento, para isso, algumas métricas como análise dos gráficos de desempenho e comparação dos valores gerados pelo modelo com os valores do mundo real, avaliando se os valores gerados pelo modelo estão apresentando uma tendência parecida com os do mundo real, dada a dificuldade de se prever os valores com exatidão, a previsão de tendência de comportamento se torna mais acessível, pois mesmo que não se chegue a um resultado 100\% preciso, uma previsão correta da tendência de comportamento também pode ser bastante útil no mercado de ações. Outro critério importante é a estabilidade do treinamento, pensando em qual modelo apresenta uma curva de aprendizado mais consistente, buscando evitar o \textit{Overfiting}.
